\section{Discussion and Limitations}
\label{sec:discussion}
\textbf{What restoration achieves.} With 16 directions at one layer, restoration reduces command error in most tasks, and the paired outcomes of the test follow a consistent pattern (Supplement, Sec.~S12.4). Failures turned into successes grow with how much of the clean state is copied: 3 to 6 under the controls, 7 under restoration, and 10 when the whole L9 state is replaced. Restoration itself gains 7 successes and loses 5 under Camera, and gains 7 and loses 3 under Sensor Noise. Every edited run, including the random-direction control, also turns 3 to 5 of the 96 Camera or 120 Sensor Noise cases from success into failure, much as a clean rerun does.

\textbf{Limitations and next steps.} Sensitivity is measured at the start of an episode; later, contact-critical steps are left to future work. Restoration is tested in $\pi_{0.5}$ at one layer, with 16 directions, on two of the four families and 216 cases, and it needs a clean observation of the same state, which deployment lacks. The shared directions also carry the response to Background, which costs $\pi_{0.5}$ no net success (39--65\% of its response energy lies along directions fitted to the other families), so they describe how the policy responds to a visual change, harmful or not. All experiments are simulated, use two related policies, and use settings chosen to keep perturbations moderate. Each limit names a next step, from tracking sensitivity through an episode to estimating clean values without a clean observation.
